% source: erdos-872/current_state.md (Shield Reduction Theorem)
% source: erdos-872/lean/README.md

\section{The Shield Reduction Theorem}\label{sec:shield-reduction}

We begin with the structural reduction that turns terminal game positions into a
weighted lower-half covering problem.  This theorem is strategy-independent and
is formally verified in Lean with no remaining sorries; see
\cite{Bloom872} for the problem statement and the repository formalization index
for the machine-checked artifact.

\subsection{Definitions}

For $x \in L$, define its upper-half multiple set by
\[
  M_n(x) := \{u \in \U : x \mid u\},
\]
and let
\[
  w_n(x) := |M_n(x)|.
\]
If $P \subseteq \U$ is a set of upper-half moves already taken by Prolonger, we
write
\[
  L(P) := \{x \in L : x \nmid u \text{ for every } u \in P\}.
\]
Thus $L(P)$ is the set of lower-half integers not yet blocked by Prolonger's
upper-half play.

\begin{definition}
For $P \subseteq \U$, define
\[
  \beta(P)
  :=
  \max\left\{
    \sum_{x \in B} w_n(x) :
    B \subseteq L(P) \text{ is an antichain}
  \right\}.
\]
\end{definition}

The quantity $\beta(P)$ is the largest weighted lower-half antichain that can
survive after Prolonger has played the shield set $P$.

\subsection{The reduction}

% source: erdos-872/current_state.md (Shield Reduction Theorem)
\begin{theorem}[Shield Reduction]\label{thm:shield-reduction}
Let $A$ be any terminal position in the divisibility antichain game, and let
$P := A \cap \U$ be the subset of upper-half moves played by Prolonger.  Then
\[
  |A| \ge |\U| - \beta(P).
\]
\end{theorem}

\begin{proof}
Set $B := A \cap L$.  Since $A$ is an antichain, $B$ is an antichain in the
lower half.  Moreover, if $x \in B$ and $u \in P$, then $x \nmid u$, because
otherwise $x$ and $u$ would be comparable in $A$.  Hence $B \subseteq L(P)$.

Now let $u \in \U \setminus A$.  Since $A$ is terminal, the move $u$ is not
legal.  Because distinct elements of $\U$ are incomparable, the only possible
reason is that some $x \in B$ divides $u$.  Therefore
\[
  \U \setminus A \subseteq \bigcup_{x \in B} M_n(x).
\]
Taking cardinalities and using the union bound yields
\[
  |\U| - |A \cap \U|
  \le
  \sum_{x \in B} |M_n(x)|
  =
  \sum_{x \in B} w_n(x)
  \le
  \beta(P).
\]
Finally,
\[
  |A|
  =
  |A \cap \U| + |A \cap L|
  \ge
  |A \cap \U|
  \ge
  |\U| - \beta(P),
\]
as claimed.
\end{proof}

\begin{remark}
The proof uses only two facts: first, every subset of $\U$ is an antichain;
second, terminality in the game means that every unused upper-half move is
blocked by some chosen lower-half divisor.  In particular, the theorem is not
tied to any particular strategy.
\end{remark}

\subsection{Why the theorem matters}

The theorem isolates the one weighted parameter that a Prolonger lower-bound
strategy must control.  If Prolonger can arrange a shield set $P$ with
\[
  \beta(P) \le cn
\]
for some fixed $c<1/2$, then
\[
  |A| \ge \left(\frac12-c+o(1)\right)n,
\]
so the game has linear length.  Conversely, if every plausible shield family
keeps $\beta(P)$ above $n/2$, then the reduction is vacuous for linear lower
bounds.

This is exactly what Theorem~A will quantify in the next section: polynomially
sized shield prefixes already face a sharp $\beta(P)$ barrier.  The point of
\Cref{thm:shield-reduction} is therefore twofold.  It is a clean theorem in its
own right, and it explains why the later obstruction results are genuinely
about the geometry of weighted primitive sets rather than about game-tree
pathologies.
